\section{Denominators in the canonical bundle formula}\label{sec:fibration}

The curve weight theorem controls a coefficient of the moduli divisor after
restriction to a transverse curve.  To use that control for a linear system,
we must keep track of the actual divisors in the canonical bundle formula,
including their principal parts.  We first choose a presentation with a
uniform pluricanonical degree, and then show that its moduli divisor has a
uniform denominator.  Throughout this section the ground field is $\C$.

\begin{proposition}[An exact presentation]\label{prop:exact-presentation}
Fix an integer $d\geq 1$ and a finite set
$\Phi\subset [0,1]\cap\Q$.  There is an integer $p_0=p_0(d,\Phi)>0$,
clearing the denominators of $\Phi$, with the following property.
Let $f:X\to Z$ be a contraction of normal projective varieties with
$\dim X\leq d$, and let $(X,B)$ be lc with $B\geq0$ and nonzero
coefficients in $\Phi$.  Suppose that $K_X+B$ is $\Q$-Cartier and
\[
 K_X+B\simq f^*L
\]
for a $\Q$-Cartier divisor $L$ on $Z$.  Then there are
$\psi\in\C(X)^*$ and a $\Q$-Cartier divisor $D_Z\simq L$ such that
\begin{equation}\label{eq:exact-presentation}
 K_X+B+\frac{1}{p_0}\Div(\psi)=f^*D_Z.
\end{equation}
\end{proposition}

\begin{proof}
Put $K=\C(Z)$ and let $F=X\times_Z\Spec K$ be the generic fibre.
Since $f$ is a contraction and the characteristic is zero, $F$ is
geometrically integral.  It is normal, its induced pair $(F,B_F)$ is
geometrically lc, and $K_F+B_F\simq0$.  These assertions may be checked
on a log resolution: after shrinking the base, the resolution and its
marked strata have the required generic smoothness, and the discrepancy
formula restricts to the fibres.  We choose the canonical divisor on $F$
by dividing a rational top form on $X$ by a rational top form on $Z$.

Theorem~\ref{thm:normal-index}, applied to the geometric generic fibre,
gives a principal multiple in a degree depending only on its dimension
and $\Phi$.  Taking a common multiple over dimensions at most $d$, and
also clearing $\Phi$, gives $p_0$.  The theorem applies over the algebraic
closure of $K$ by characteristic-zero comparison: all data descend to an
algebraically closed field admitting an embedding into $\C$.

By Lemma~\ref{lem:trivialization-descent}, this trivialization descends
to $K$ in the same degree.  Consequently we may choose
$\psi\in\C(X)^*$ such that
\[
 V:=p_0(K_X+B)+\Div(\psi)
\]
has no component dominating $Z$.

It remains to show that this particular vertical divisor is a pullback.
Choose $n>0$ and $h\in\C(X)^*$ with
$n(K_X+B-f^*L)=\Div(h)$.  Then
\[
 n(V-p_0f^*L)=\Div(h^{p_0}\psi^n).
\]
The divisor on the right is vertical.  Its defining function has zero
divisor on the normal projective generic fibre and hence belongs to
$K^*$, because $H^0(F,\OO_F)=K$.  Write
$h^{p_0}\psi^n=f^*a$ for $a\in K^*$.  We obtain
\[
 V=f^*\left(p_0L+\frac1n\Div(a)\right).
\]
Thus $D_Z=L+(np_0)^{-1}\Div(a)$ has all the asserted properties.
\end{proof}

We recall the part of the canonical bundle formula needed below.  Fix a
presentation~\eqref{eq:exact-presentation}.  For a prime divisor $P$ on
$Z$, let $t_P$ be the log canonical threshold of $f^*P$ over the generic
point of $P$; here $P$ is Cartier after restricting to a neighborhood of
that point.  Set
\begin{equation}\label{eq:cbf-definition}
 B_Z=\sum_P(1-t_P)P,\qquad M_Z=D_Z-K_Z-B_Z.
\end{equation}
On a higher model $\tau:W\to Z$, use the crepant induced sub-pair to
define the thresholds and put
\begin{equation}\label{eq:cbf-higher}
 D_W=\tau^*D_Z=K_W+B_W+M_W,
\end{equation}
with compatible canonical divisors.  A sub-pair here permits negative
boundary coefficients.  The traces $M_W$ form the moduli b-divisor
$\mathbf M$.  We say that it is \emph{determined on $W$} if its trace on
every higher model is the pullback of $M_W$.

The lc-trivial fibration theorem gives a projective model on which
$\mathbf M$ is determined and its trace is nef and $\Q$-Cartier
\cite[Theorem~3.6]{FG}, extending the klt-trivial theory of
Ambro~\cite{AmbroBoundary,Ambro}.  The lc formulation
uses the discrepancy b-divisor $\mathbf A^*$ with the discrepancy $-1$
terms omitted.  Its rank-one hypothesis holds here: on a resolution,
$\lceil\mathbf A^*\rceil$ has only effective exceptional terms over the
generic fibre, since $B$ is an effective boundary.  Its pushforward on
that normal fibre is the structure sheaf, and the contraction condition
gives rank one.  In particular, horizontal components of coefficient one
are allowed.  A further resolution gives a smooth projective
determination.  We use the representative~\eqref{eq:cbf-definition};
changing $D_Z$ by a rational principal divisor changes $\mathbf M$ by
the corresponding principal b-divisor and preserves these qualitative
properties.

The original discriminant $B_Z$ is effective and its coefficients lie in
a rational DCC set depending only on $d$ and $\Phi$.  To see this, lc
gives $t_P\geq0$.  A component of $f^*P$ with multiplicity $a\geq1$
and boundary coefficient $b\geq0$ gives
$t_P\leq(1-b)/a\leq1$.  Near the generic point of $P$, the testing
divisor $f^*P$ has positive integral coefficients.  A log resolution,
followed by removing proper closed subsets of $P$, realizes its generic
threshold as an ordinary threshold on an open subset of $X$.
The ACC theorem for log canonical thresholds
\cite[Theorem~1.1]{HMX}, with testing coefficient set
$\Z_{>0}$, therefore puts the numbers $1-t_P$ in the claimed DCC set.
On higher models, $B_W$ need not be effective, but its coefficients
remain at most one, since the induced sub-pair is crepant and sub-lc.

\begin{theorem}[A uniform denominator for the moduli divisor]
\label{thm:moduli-denominator}
Fix $d\geq1$ and a finite set $\Phi\subset[0,1]\cap\Q$, and let
$p_0=p_0(d,\Phi)$ be as in
Proposition~\ref{prop:exact-presentation}.  There is a positive integer
$p=p(d,\Phi)$ divisible by $p_0$ with the following property.
Let $f:X\to Z$ be a contraction of normal projective complex varieties
with $\dim X\leq d$ and $\dim Z>0$, and let $(X,B)$ be lc with
$B\geq0$, coefficients in $\Phi$, and $K_X+B$ $\Q$-Cartier.
For any exact presentation
\[
 K_X+B+\frac1{p_0}\Div(\psi)=f^*D_Z
\]
with $\psi\in\C(X)^*$ and $D_Z$ $\Q$-Cartier, let $\mathbf M$ be
its moduli b-divisor.
On every smooth projective determination $W$ of $\mathbf M$,
the divisor $pM_W$ is nef Cartier.
\end{theorem}

The normalization by $p_0$ fixes the scale at which principal changes
are allowed.  The next lemma computes one coefficient of this actual
moduli divisor by a curve weight.  It does not require the crepant
boundary on a resolution to be effective.

\begin{lemma}[A transverse curve and its weight]\label{lem:slice}
Let $f:X\to Z$ be a contraction of normal projective complex varieties,
let $(X,B)$ be lc with $B\geq0$, and suppose that
\[
 K_X+B+\frac1m\Div(\psi)=f^*D_Z
\]
for an integer $m>0$, a function $\psi\in\C(X)^*$, and a
$\Q$-Cartier divisor $D_Z$.
Let $\tau:W\to Z$ be a smooth projective birational model, and let
$P$ be a prime divisor of $W$.  Write
$\alpha=\operatorname{coeff}_P(\tau^*D_Z)$, and let $t_P$ be the
threshold for the crepant induced sub-pair over the generic point of
$P$.  Put $s=\dim X-\dim Z$.
Then there are a smooth projective curve $C$, a point $c\in C$, a
rational uniformizer $z$ at $c$, and a regular function-field extension
of $\C(C)$ admitting a geometrically integral projective normal
effective $s$-fold pair $(V,B_V)$ that is geometrically lc, and a
rational relative $m$-canonical form $\phi$ such that
\begin{equation}\label{eq:slice-weight}
 \Div(\phi)+mB_V=0,\qquad w_z(\phi)=\alpha-1+t_P.
\end{equation}
\end{lemma}

\begin{proof}
Choose a smooth projective common resolution $U$ mapping to $X$ and to
$W$, and denote the latter morphism by $g$.  With a compatible rational
top form $\theta$ on $U$, write
\begin{equation}\label{eq:slice-upstairs}
 \Div(\theta)+B_U^c+\frac1m\Div(\psi)=g^*D_W,
 \qquad D_W=\tau^*D_Z.
\end{equation}
Here $B_U^c$ is the crepant sub-boundary.  Resolve so that its support,
the exceptional divisors over $X$, and the support of $g^*P$ are SNC.
We may include the finitely many canonical and principal supports
occurring in~\eqref{eq:slice-upstairs} among the markings.

Write $r=\dim Z$.  If $r>1$, take a sufficiently general smooth
complete-intersection curve
\[
 C=L_1\cap\cdots\cap L_{r-1}\subset W
\]
of very ample members, and choose a transverse point $c\in C\cap P$
over a general point of $P$.  If $r=1$, take $C=W$ and $c=P$.
Here are the avoidance and smoothness conditions used in this choice.
The fibre-dimension jumping locus of $g$ has codimension at least two:
a jump over a divisor would give a proper inverse image of dimension
$\dim U$.  There are also finitely many smooth marked intersections
on $U$.  We avoid their images when those images have codimension at
least two.  For each intersection whose image is a divisor, generic
smoothness gives a dense smooth open of that divisor over which its
map is smooth; we avoid the complement.  All these excluded closed
sets have codimension at least two in $W$, so a general complete
intersection curve misses them.  Choose $c$ also away from the other
divisor supports.

For global smoothness of the cut, apply Bertini successively to the
basepoint-free systems $g^*|L_j|$ on $U$ and simultaneously on its
smooth marked intersections.  This gives a smooth variety
$\widetilde U=g^{-1}(C)$ with restricted SNC markings.  This argument
uses the smoothness of $U$, not smoothness or flatness of $g$ along
$P$.

The morphism $g$ has connected fibres, as follows from Stein
factorization and the fact that $\C(W)$ is relatively algebraically
closed in $\C(U)$.  Its restriction $h:\widetilde U\to C$ has the
same fibres over points of $C$.  Thus $\widetilde U$ is connected and,
being smooth, is integral; $h$ is a contraction.  In particular its
function field is regular over $\C(C)$.

We record why the threshold is retained in this cut.  On the SNC model,
the threshold over the generic point of $P$ is
\[
 \min_i\frac{1-b_i}{a_i},
\]
where the marked prime divisors dominating $P$ have crepant boundary
coefficients $b_i$ and multiplicities $a_i$ in $g^*P$.
Transversality gives the same coefficients and multiplicities in the
fibre over $c$.  Components with image a proper closed subset of $P$
are avoided by the choice of $c$.  The SNC criterion for sub-lc pairs
then shows that this minimum is also the threshold on the cut.

Exact adjunction fixes the coefficient of the base divisor as well.
For each $j$ choose $L'_j\sim L_j$ avoiding $c$, and choose a rational
function on $W$ with divisor $L'_j-L_j$.  Multiply $\theta$ by the
pullbacks of these functions and take iterated residues along the cuts.
The resulting rational top form $\widetilde\theta$ on
$\widetilde U$ has the canonical divisor prescribed by adjunction,
namely the restriction of
$\Div(\theta)+\sum_jg^*L'_j$.  Restricting
\eqref{eq:slice-upstairs} therefore gives
\begin{equation}\label{eq:slice-exact}
 \Div(\widetilde\theta)+\widetilde B^c+
 \frac1m\Div(\widetilde\psi)=h^*D'_C,
 \qquad
 D'_C=\left(D_W+\sum_j L'_j\right)\big|_C.
\end{equation}
All restrictions are defined by generality.  Since $L'_j$ avoid $c$
and $C$ meets $P$ transversely,
$\operatorname{coeff}_cD'_C=\alpha$.  When $r=1$, the sums and the
residue operations are empty.

To obtain the effective generic presentation, return to the original
model $X$.  Its normal generic fibre over $Z$ is geometrically normal
over its perfect characteristic-zero ground field, and is geometrically
integral since $f$ is a contraction.  Generic flatness and openness of
geometric normality and integrality therefore give a dense open subset
of $W$ where $W\to Z$ is an isomorphism and the original fibres have
these properties.  Shrink this open set further using generic smoothness
of the resolution strata, so that the restricted resolution computes
the fibre discrepancies also after algebraic closure.
Exceptional centres that do not dominate the base disappear
there; centres that dominate it retain codimension at least two in
the fibres by the dimension formula.  These are conditions at the
generic point of $C$, which is chosen in this open set even when the
marked point $c$ is outside it.  Hence the generic
fibre $V$ is a normal projective geometrically lc pair with the
effective boundary $B_V$ induced by $B$, and $\widetilde B^c$ compares
crepantly with it.

Choose a rational uniformizer $z$ at $c$ and define the relative form
\[
 \phi=(\widetilde\theta/dz)^{\otimes m}\widetilde\psi.
\]
Restriction of~\eqref{eq:slice-exact} to the generic fibre gives
$\Div(\phi)+mB_V=0$.  The absolute form used to compute its weight is
\[
 \Omega=\phi\wedge(dz/z)^{\otimes m}
       =z^{-m}\widetilde\theta^{\otimes m}\widetilde\psi.
\]
For a divisor $E$ over the cut lying above $c$, put
$a_E=\ord_E(z)>0$ and let $b_E$ be its crepant boundary coefficient.
Equation~\eqref{eq:slice-exact} gives
\[
 \frac1m\ord_E(\Omega)+1
       =(\alpha-1)a_E+(1-b_E).
\]
Taking the infimum after division by $a_E$ proves
\eqref{eq:slice-weight}, since the infimum of
$(1-b_E)/a_E$ is exactly the preserved threshold $t_P$.
\end{proof}

\begin{proof}[Proof of Theorem~\ref{thm:moduli-denominator}]
Let $P$ be any prime divisor of the smooth determination $W$, and put
$\alpha=\operatorname{coeff}_P D_W$.  From
\eqref{eq:cbf-higher},
\[
 \operatorname{coeff}_P M_W
   =\alpha+t_P-1-\operatorname{coeff}_P K_W.
\]
Lemma~\ref{lem:slice} realizes $\alpha-1+t_P$ as the weight of an
exact degree-$p_0$ form with an effective lc generic presentation of
dimension $s=\dim X-\dim Z$.  Theorem~\ref{thm:weight} consequently
gives
\[
 q(s,p_0)\operatorname{coeff}_P M_W\in\Z,
\]
because $K_W$ is integral.  Take $p$ to be a common multiple of $p_0$
and $q(s,p_0)$ for $0\leq s\leq d-1$.  This choice works for every
$P$ without any bound on the complexity of $W$.  Thus $pM_W$ is an
integral Weil divisor on the smooth variety $W$, hence Cartier.
It is nef by the lc-trivial fibration theorem.
\end{proof}
